\section{The decay length of Hastings' approximate factorization}
\label{sec:hastings_decay_slack}

% Hastings, arXiv:0705.2024v4, law16.tex, lines 903--907 (the final bound of
% the proof of Lemma 1) and lines 141--142 (the printed decay length).

The final estimate of the proof of Lemma~1 of \cite{Hastings2007AreaLaw},
lines 903--907, bounds the error of the approximate factorization by a
multiple of $\sqrt{l/\xi}\,e^{-l/6\xi}$. The prefactor $\sqrt l$ cannot be
absorbed into a constant at the rate $e^{-l/6\xi}$, but it can at every
slower rate $e^{-l/\xi'}$ with $\xi'>6\xi$.

\begin{theorem}[A square root against an exponential]
    \label{thm:hastings_sqrt_exp}
    \lean{TNLean.MPS.AreaLaw.Hastings.sqrt_mul_exp_neg_le}
    \leanok
    For $\delta>0$ and $l\ge0$,
    \begin{align}
        \sqrt l\,e^{-\delta l} & \le\frac1{\sqrt{2e\delta}},
        \label{eq:hastings_sqrt_exp}
    \end{align}
\end{theorem}
\begin{proof}\leanok
    With $x=2\delta l$, the inequality $x\le e^{x-1}$ gives $xe^{-x}\le e^{-1}$,
    that is $2e\delta\,l\,e^{-2\delta l}\le1$, which is the square
    of~\eqref{eq:hastings_sqrt_exp}. The bound is attained at $l=1/(2\delta)$.
\end{proof}

\begin{theorem}[A power against an exponential]
    \label{thm:hastings_pow_exp}
    \lean{TNLean.MPS.AreaLaw.Hastings.pow_mul_exp_neg_le,
        TNLean.MPS.AreaLaw.Hastings.pow_mul_exp_neg_add_le}
    \leanok
    For $\delta>0$, a natural number $p$, a real number $r$ and $l\ge0$,
    \begin{align}
        l^pe^{-\delta l}    & \le\Bigl(\frac{p}{e\delta}\Bigr)^p,        &
        l^pe^{-(r+\delta)l} & \le\Bigl(\frac{p}{e\delta}\Bigr)^pe^{-rl}.
        \label{eq:hastings_pow_exp}
    \end{align}
\end{theorem}
\begin{proof}\leanok
    With $x=\delta l/p$ for $p\ge1$, $l^pe^{-\delta l}=(p/\delta)^p(xe^{-x})^p$ and
    $xe^{-x}\le e^{-1}$. The second bound follows from the first after
    splitting off $e^{-rl}$.
\end{proof}

\begin{theorem}[Slack in the decay length]
    \label{thm:hastings_decay_slack}
    \lean{TNLean.MPS.AreaLaw.Hastings.add_mul_sqrt_mul_exp_le_of_lt}
    \leanok
    Let $\xi>0$, $\xi'>6\xi$ and $\delta=1/(6\xi)-1/\xi'$. For $a,b\ge0$ and
    $l\ge0$,
    \begin{align}
        (a+b\sqrt l)\,e^{-l/6\xi} & \le\Bigl(a+\frac{b}{\sqrt{2e\delta}}\Bigr)e^{-l/\xi'}.
        \label{eq:hastings_decay_slack}
    \end{align}
    Applied to the final estimate of the proof of Lemma~1, this gives the
    approximate factorization with error $C_1e^{-l/\xi'}$ for every
    $\xi'>6\xi$, with $C_1$ depending on $\xi'-6\xi$, rather than at the
    printed $\xi'=6\xi$ of \cite[lines 141--142]{Hastings2007AreaLaw}.
\end{theorem}
\begin{proof}\leanok
    \uses{thm:hastings_sqrt_exp}
    Write $e^{-l/6\xi}=e^{-l/\xi'}e^{-\delta l}$. Then $a\,e^{-\delta l}\le a$,
    and $b\sqrt l\,e^{-\delta l}\le b/\sqrt{2e\delta}$
    by~\eqref{eq:hastings_sqrt_exp}.
\end{proof}

\begin{theorem}[From the square-root bound to a decay length beyond $6\xi$]
    \label{thm:hastings_factorization_rate}
    \lean{TNLean.MPS.AreaLaw.Hastings.exists_isFactorizationAt_of_sqrt_decay}
    \leanok
    \uses{def:hastings_factorization_at}
    Let $\xi>0$, $\xi'>6\xi$, $\delta=1/(6\xi)-1/\xi'$ and $a,b\ge0$. Suppose
    that at a cut $c$ and length $l$ there is an approximate factorization
    with error $(a+b\sqrt l)\,e^{-l/6\xi}$. Then at the same $(c,l)$ there is
    one with error
    \begin{align}
        \Bigl(a+\frac{b}{\sqrt{2e\delta}}\Bigr)e^{-l/\xi'}.
        \label{eq:hastings_factorization_rate}
    \end{align}
    This is a stage of Lemma~1 of \cite{Hastings2007AreaLaw}: the
    factorizations with the square-root error, the form the proof of Lemma~1
    reaches at lines 903--907, are a hypothesis, and the theorem only converts
    their error to the form $C_1e^{-l/\xi'}$ of the conclusion of Lemma~1, for
    a decay length $\xi'>6\xi$ in place of the printed
    $\xi'=6\xi$~\cite{gap:hastings07_gaussian_average_exponent}.
\end{theorem}
\begin{proof}\leanok
    \uses{thm:hastings_decay_slack, thm:hastings_factorization_mono}
    By Theorem~\ref{thm:hastings_decay_slack} the error
    $(a+b\sqrt l)\,e^{-l/6\xi}$ is at most~\eqref{eq:hastings_factorization_rate},
    and a factorization with a smaller error is one with a larger error.
\end{proof}

\begin{theorem}[The positive form at a decay length beyond $6\xi$]
    \label{thm:hastings_positive_factorization_rate}
    \lean{TNLean.MPS.AreaLaw.Hastings.exists_posSemidef_isFactorizationAt_of_sqrt_decay}
    \leanok
    \uses{def:hastings_factorization_at}
    Under the hypotheses of Theorem~\ref{thm:hastings_factorization_rate},
    with $\Omega$ a unit vector and $C_1=a+b/\sqrt{2e\delta}$, there is at the
    same $(c,l)$ an approximate factorization with $O_B$ positive
    semidefinite and error $(\sqrt2+4)\sqrt{C_1}\,e^{-l/(2\xi')}$.
\end{theorem}
\begin{proof}\leanok
    \uses{thm:hastings_factorization_rate, thm:hastings_positive_factorization}
    Apply Theorem~\ref{thm:hastings_positive_factorization} to the
    factorizations of Theorem~\ref{thm:hastings_factorization_rate}.
\end{proof}
